\chapter*{RÉSUMÉ}
\addcontentsline{toc}{chapter}{RÉSUMÉ}
Ce mémoire de master à visée de recherche explore l'optimisation de la \textbf{reconstruction d'images issue de la tomodensitométrie} (TDM) pour concilier qualité diagnostique et réduction de la dose d'irradiation. Face aux limites de la rétroprojection filtrée (FBP) en situation de sous-échantillonnage, ce travail adopte le paradigme du Compressive Sensing en couplant l'algorithme d'optimisation \textbf{ADMM} aux ondelettes biorthogonales \textbf{Bior4.4}. Ce choix technique permet d'exploiter la parcimonie des images médicales pour reconstruire fidèlement les structures fines à partir de mesures réduites.
Structuré en quatre chapitres allant des fondements théoriques (Radon, Fourier) à l'expérimentation numérique, ce mémoire démontre la supériorité de l'approche proposée. Sur 19 images tests, la méthode atteint un PSNR moyen de $26.20\,dB$ et un SSIM de $0,866$, surpassant nettement la FBP (21,61 dB) et les algorithmes ISTA/FISTA. Avec une convergence stable et un taux de parcimonie de 20\%, ces résultats ouvrent des perspectives majeures : accélération sur GPU pour le temps réel, passage à la reconstruction 3D et hybridation avec l'apprentissage profond.